\section{The F Track: Paper Overview}
\label{sec:overview}

The F track comprises six companion papers. This section briefly previews
each one.

\subsection{F.1 --- Single-Chamber All 50 States}
\label{sec:overview:f1}

\textbf{Title}: ``Algorithmic State House Redistricting: A 50-State Empirical Study''

F.1 is the primary empirical paper. It applies the \texttt{redist} pipeline
with adaptive resolution selection (\texttt{block\_group} when
$k/n_{\mathrm{tracts}} > 0.05$, \texttt{tract} otherwise) to all 50 state
house chambers. The paper reports Polsby-Popper scores, population balance
statistics, runtime measurements, and county-split counts for every state.
Four case studies are developed in depth: Washington (98 districts),
Texas (150), New Hampshire (400), and Nebraska's unicameral legislature (49).

\subsection{F.2 --- Bicameral Nesting}
\label{sec:overview:f2}

\textbf{Title}: ``NestSection at Scale: Spine-Compatible Bicameral Redistricting for All 50 States''

F.2 applies the NestSection algorithm \citep{dellaLibera2026nestSection} to
all 49 bicameral state legislatures. NestSection constructs a geographic spine
at the level of $\gcd(H,S)$ sub-regions from which both house and senate maps
are derived. F.2 classifies all 50 states by the compatibility of their
$H:S$ ratio, estimates the compactness penalty from nesting, and measures
the partisan effect of requiring senate districts to be exact unions of
house districts.

\subsection{F.3 --- Multi-Resolution High-$k$ Chambers}
\label{sec:overview:f3}

\textbf{Title}: ``Resolution Selection for State Legislative Redistricting: When Census Tracts Are Too Coarse''

F.3 is the methods paper for resolution selection. It derives the
$k/n_{\mathrm{tracts}} > 0.05$ threshold theoretically and validates it
empirically on three high-$k$ states (Washington, Texas, California).
The paper quantifies the MAUP effects of resolution choice on partisan
and compactness outcomes, and supplies a resolution recommendation table
for all 50 states and all lower chambers.

\subsection{F.4 --- State-Level Criteria Variation}
\label{sec:overview:f4}

\textbf{Title}: ``Algorithmic Redistricting Under Heterogeneous State Criteria''

Many states impose explicit compactness requirements (e.g., Ohio's
``no more splitting of political subdivisions than necessary'' standard),
municipal-boundary preservation rules, or community-of-interest requirements
that go beyond the federal one-person-one-vote floor.
F.4 surveys these state-specific criteria and evaluates how the \texttt{redist}
pipeline's objective function can be adapted to meet them.

\subsection{F.5 --- Compactness: Legislative vs.\ Congressional}
\label{sec:overview:f5}

\textbf{Title}: ``Why Are State Legislative Districts More Compact? A Scale Analysis''

The headline finding of F.1---that state legislative maps achieve higher
Polsby-Popper scores than congressional maps---is counterintuitive.
Smaller districts might be expected to be \emph{less} compact because they
are more sensitive to local population concentrations.
F.5 resolves this paradox by decomposing PP into a scale component (smaller
perimeter for smaller districts) and a shape component (algorithm performance
at different $k$ values). The analysis shows that the scale component
dominates: PP is primarily determined by the ratio of district area to
perimeter squared, and smaller districts have smaller perimeters relative
to their area when drawn compactly.

\subsection{F.6 --- VRA Compliance at the State Level}
\label{sec:overview:f6}

\textbf{Title}: ``Voting Rights Act Compliance in Algorithmic State Legislative Redistricting''

F.6 extends the D-track VRA analysis \citep{dellaLibera2026vraCompliance}
to state legislative chambers. State-level minority opportunity district
requirements arise from both the federal VRA and state constitutional
provisions (e.g., California's Fair Maps Act). F.6 applies the
VRA-Section algorithm with edge-weighting calibrated to minority
population concentrations at the block-group level.
